\subsection{诺特环上的有界平坦复形}

命题11。——设 A 是左诺特环，令 C 为平坦左 A-模的有界复形，且 H(C) 是有限生成 A-模。设 a 和 b 为两个整数，使得 \(a \leqslant b\) 并且 \(\mathrm{H}_{n}(\mathrm{C}) = 0\) （当 \(n < a\) ）， \(\mathrm{C}_{n} = 0\) （当 \(n > b\) ）。存在一个左 A-模复形 P，使得 \(\mathrm{P}_{n}\) 对每个 \(n\) 都是投射且有限生成的，并且 \(\mathrm{P}_{n} = 0\) （当 \(n \notin [a, b]\) ），以及一个拟同构 \(u: \mathrm{P} \to \mathrm{C}\) 。此外，对于右 A-模的每个复形 E，同态

\[
\mathrm{H} \left(1 _ {\mathrm{E}} \otimes u\right): \mathrm{H} (\mathrm{E} \otimes_ {\mathrm{A}} \mathrm{P}) \rightarrow \mathrm{H} (\mathrm{E} \otimes_ {\mathrm{A}} \mathrm{C}) \quad \text{是双射。}
\]

根据 X，第 53 页，命题 7，存在一个复形 (L, d)，使得 \(L_{n}\) 对每个 n 都是自由且有限生成的，并且当 n \textless{} a 时为零，以及一个拟同构 \(f: L \to C\) 。设 P 为商复形 \(L/L'\) ，其中 \(L_{n}' = 0\) （当 n \textless{} b ）， \(L_{n}' = L_{n}\) （当 n \textgreater{} b ），以及 \(L_{b}' = B_{b}(L)\) 。由于 \(C_{n} = 0\) （当 n \textgreater{} b ），故 \(f(L') = 0\) ，因此 f 通过一个复形态射 \(u: P \to C\) 分解.

\[
\begin{array}{c c c c c} \dots \longrightarrow & L _ {b + 1} & \xrightarrow {d _ {b + 1}} & L _ {b} & \xrightarrow {d _ {b}} L _ {b - 1} \longrightarrow \dots \\ & \downarrow & & \downarrow & | | \\ & 0 & \longrightarrow & P _ {b} & \longrightarrow P _ {b - 1} \longrightarrow \dots \\ & \downarrow & & \downarrow^ {u _ {b}} & \downarrow^ {u _ {b - 1}} \\ & 0 & \longrightarrow & C _ {b} & \longrightarrow C _ {b - 1} \longrightarrow \dots \end{array}
\]

由于 \(f\) 是一个拟同构，我们有 \(\mathrm{H}(\mathrm{Con}(f)) = 0\) ，由此得到正合序列

\[
\dots \rightarrow \mathrm{L} _ {b + 1} \xrightarrow {d _ {b + 1}} \mathrm{L} _ {b} \longrightarrow \mathrm{L} _ {b - 1} \oplus \mathrm{C} _ {b} \longrightarrow \mathrm{L} _ {b - 2} \oplus \mathrm{C} _ {b - 1} \rightarrow \dots .
\]

因此我们有一个正合序列

\[
0 \rightarrow \mathrm{P} _ {b} \rightarrow \mathrm{L} _ {b - 1} \oplus \mathrm{C} _ {b} \rightarrow \mathrm{L} _ {b - 2} \oplus \mathrm{C} _ {b - 1} \rightarrow \dots .
\]

这表明，一方面， \(u\) 的锥体具有零同调，因此 \(u\) 是拟同构，另一方面，模 \(\mathbf{P}_{b}\) 是平坦的（X，第 76 页，推论2）；由于 \(\mathbf{P}_{b}\) 作为 \(\mathbf{L}_{b+1}\) 的商是有限型的，它是投射的（X，第 13 页，推论）。因此，数对 (P, u) 满足所需条件。最后一个断言由 X，第 79 页，推论 5 得出。

* 例。--- 设 A 为诺特交换环，X 为 A 上的真且平坦的概形， \(\mathcal{F}\) 为一个凝聚 \(\mathcal{O}_{\mathbf{X}}\)-模，在 A 上平坦。存在一个有界复形 P，由有限生成的投射 A-模组成，使得对每个 A-模 M，H(X, \(\mathcal{F} \otimes_{\mathbf{A}} \mathbf{M}\) ) 自然等同于 H(P \(\otimes_{\mathbf{A}} \mathbf{M}\) ) 。事实上，设 \(\mathfrak{U}\) 为 X 的由有限个仿射开集组成的覆盖， \(\mathcal{E}(\mathfrak{U}, \mathcal{F})\) 为相应的 Čech 复形：可证明 H \(^{i}\) ( \(\mathcal{E}(\mathfrak{U}, \mathcal{F})\) ) 同构于 A-模 H \(^{i}\) (X, \(\mathcal{F}\) ) ，且后者是有限生成的；此外，对每个 A-模 M，复形 \(\mathcal{E}(\mathfrak{U}, \mathcal{F}) \otimes_{\mathbf{A}} \mathbf{M}\) 同构于 \(\mathcal{E}(\mathfrak{U}, \mathcal{F} \otimes_{\mathbf{A}} \mathbf{M})\) 。将命题11应用于复形 \(\mathcal{E}(\mathfrak{U}, \mathcal{F})\) （它是有界的），得到一个满足要求的复形 P。

设 \(y\) 为 \(\text{Spec}(\mathbf{A})\) 的一点， \(\kappa(y)\) 为 \(\mathbf{A}\) 在 \(y\) 的剩余域， \(X_y = X \otimes_{\mathbf{A}} \kappa(y)\) 为 \(\mathbf{X}\) 在 \(y\) 之上的纤维， \(\mathcal{F}_y = \mathcal{F} \otimes_{\mathbf{A}} \kappa(y)\) ，并令 \(h^p(y) = \dim_{\kappa(y)} H^p(X_y, \mathcal{F}_y)\) （当 \(p \geqslant 0\) ）。由复形 \(\mathbf{P}\) 的存在性容易推出以下结果：

\begin{enumerate}
\def\labelenumi{(\roman{enumi})}
\item
  函数 \(h^{p}\) 在 Spec(A) 上是上半连续的；
\item
  函数 \(\sum_{p\geqslant 0}(-1)^{p}h^{p}\) 在 \(\operatorname {Spec}(\mathbf{A})\) 上是局部常数的.*
\end{enumerate}

